\begin{center}
    \underline{\textbf{AUDIT INSTRUMENTASI METROLOGIS: KETEPATAN DAN}} \\
    \underline{\textbf{KETIDAKPASTIAN PADA EKSPERIMEN P1 (\textit{DENSITY}) DAN M3 (\textit{ENERGY})}}
\end{center}

\begin{center}
    Muhammad Abiyu Khansa (182241036), Farras Atha Pandita (182241040), Hendra Maulana (182241043)
\end{center}

\hrule height 2pt

\begin{center}
    \textbf{Program Studi Fisika, Departemen Fisika, Fakultas Sains dan Teknologi, Universitas Airlangga, Surabaya}
\end{center}

\vspace{10pt}
\renewcommand{\abstractname}{\normalsize{ABSTRAK}}
\begin{abstract}
\noindent
Laporan ini memuat audit instrumentasi metrologis Group I pada dua eksperimen, yaitu P1 (massa jenis silinder kuningan) dan M3 (konservasi energi mekanik pada rel udara). Sesuai ketentuan kelompok beranggotakan tiga orang, setiap besaran diukur sebanyak $n=8$ kali. Penapisan pencilan dengan kriteria $3\sigma$ tidak membuang satu pun data. Ketidakpastian dievaluasi menurut ISO-GUM (Tipe A dan Tipe B, $\uc$, dan $U$ pada $k=2$). Hasil P1 adalah $\rho=(8{,}493\pm0{,}035)\ \si{g/cm^3}$ dengan kontribusi ketidakpastian terbesar berasal dari diameter ($85\%$). Kalibrasi neraca digital terhadap $8$ titik massa standar menghasilkan $y=0{,}999975\,x+0{,}0421$ dengan $R^2\approx 1$, tetapi plot residual menunjukkan pola lengkung sistematis (simpangan maksimum $0{,}032\ \si{g}$) sehingga $R^2$ saja tidak cukup sebagai bukti linearitas. Pada M3 diperoleh $\Delta E=E_1-E_2=(-2{,}79\pm0{,}61)\times10^{-3}\ \si{J}$; tanda negatif berarti energi akhir lebih besar daripada energi awal, yang secara fisis tidak mungkin pada sistem disipatif. Hasil ini ditafsirkan sebagai indikasi galat sistematis (terutama \textit{response time delay} sensor gerbang cahaya) dan bukan sebagai pelanggaran hukum kekekalan energi.
\end{abstract}

\vspace{5pt}
\noindent\textbf{Kata kunci:} audit instrumentasi, ISO-GUM, ketidakpastian bentangan, kalibrasi OLS, analisis residual, massa jenis, kekekalan energi, rel udara.
